\addchap{Wstęp}

\section*{O czym jest ten kurs}

Celem kursu jest zbudowanie od zera procesora RISC-V (RV32I) w SystemVerilogu,
najpierw w wersji jednocyklowej, a potem jako klasyczny 5-etapowy potok,
i uruchomienie na nim programów skompilowanych z C. Po drodze powstaje
wszystko, czego potrzebuje projektant procesora: emulator referencyjny,
infrastruktura testowa, co-symulacja i fuzzing.

Notatki zakładają, że znasz C i programowanie niskopoziomowe (wskaźniki,
pamięć, ABI, asembler w zarysie). Nie zakładają znajomości języków opisu
sprzętu ani elektroniki. Tam, gdzie intuicja programisty C prowadzi na
manowce, znajdziesz czerwone ramki \emph{Pułapka}.

\section*{Jak czytać te notatki}

Każdy rozdział odpowiada jednemu modułowi kursu (katalogowi \plik{NN-nazwa/}).
Rozdział zaczyna się od metryczki (cel, czas, katalog), potem jest teoria z
przykładami, zadania (zielone ramki) i pytania kontrolne. Rozwiązania zadań
piszesz w plikach \plik{rtl/*.sv} danego modułu, a gotowe testy w
\plik{tb/} mówią, czy działa.

\begin{center}
\small
\begin{tabularx}{\linewidth}{@{}c l X r@{}}
\toprule
\textbf{Rozdz.} & \textbf{Moduł} & \textbf{Co budujesz} & \textbf{Wieczory} \\
\midrule
0 & \plik{00-start} & narzędzia, licznik, pierwszy testbench & 1 \\
1 & \plik{01-kombinacyjna} & mux, enkoder, sumator, barrel shifter; szerokości, zatrzaski & 2 \\
2 & \plik{02-sekwencyjna} & rejestry, timing, pamięci, FSM, FIFO, UART & 2--3 \\
3 & \plik{03-alu-regfile} & ALU, bank rejestrów, komparator skoków & 1--2 \\
4 & \plik{04-isa-emulator} & RV32I na poziomie bitów, generator immediate'ów, emulator w C & 2--3 \\
5 & \plik{05-single-cycle} & dekoder, LSU, procesor jednocyklowy & 3--4 \\
6 & \plik{06-weryfikacja} & co-symulacja, fuzzing, Verilator, mutacje & 1--2 \\
7 & \plik{07-c-bare-metal} & crt0, skrypt linkera, \code{\_\_mulsi3}: C na Twoim CPU & 1--2 \\
8 & \plik{08-pipeline} & potok 5-etapowy: forwarding, stalle, flush & 4--6 \\
9 & \plik{09-dalej} & rozszerzenie M, przerwania, FPGA, cache, formalna & --- \\
\bottomrule
\end{tabularx}
\end{center}

Moduły trzeba robić po kolei. Moduł 5 używa Twoich plików z modułów 3 i 4,
moduł 6 używa emulatora z modułu 4, a moduł 8 używa wszystkiego.

\section*{Narzędzia}

\begin{tabularx}{\linewidth}{@{}l X@{}}
\toprule
\textbf{Narzędzie} & \textbf{Do czego} \\
\midrule
Icarus Verilog (\code{iverilog}, \code{vvp}) & symulacja zdarzeniowa, 4 wartości logiczne, główny symulator kursu \\
GTKWave & oglądanie przebiegów (pliki \code{.vcd}) \\
Verilator & lint (\emph{-Wall dla sprzętu}) i szybka symulacja \\
Yosys & synteza: ile bramek, jak długa ścieżka krytyczna \\
\code{rvtool.py} & asembler i disasembler RV32I(M) dołączony do kursu \\
clang + \code{ld.lld} albo GCC dla RISC-V & tylko w rozdziale 7 (C na procesorze) \\
\bottomrule
\end{tabularx}

\medskip
Do rozdziału 7 potrzebny jest linker: \code{sudo apt install lld} (albo
pełny toolchain \code{gcc-riscv64-unknown-elf}). Pozostałe rozdziały nie
wymagają żadnego toolchaina RISC-V.

\section*{Rytm pracy}

W każdym module te same polecenia (uruchamiane z katalogu modułu):

\begin{sh}
make test               # wszystkie testy modułu (na starcie: same FAIL)
make unit T=mux4        # jeden test jednostkowy
make wave T=mux4        # przebiegi w GTKWave
make lint               # Verilator --lint-only
make synth T=adder      # Yosys: bramki i najdłuższa ścieżka
\end{sh}

W modułach z procesorem (5--8) dochodzą:

\begin{sh}
make progs              # 12 programów testowych na Twoim rdzeniu
make prog P=08_fib      # jeden program: wyjście + ślad wykonania
make cosim              # porównanie śladu z emulatorem
make fuzz N=100         # losowe programy + porównanie z emulatorem
make progs SIM=verilator
\end{sh}

\begin{uwaga}[Rada]
Zrób z katalogu kursu repozytorium git i commituj po każdym zadaniu. Przy
potoku (rozdział 8) możliwość powrotu do ostatniej działającej wersji jest
bezcenna.
\end{uwaga}

\section*{Konwencje}

\begin{itemize}
  \item SystemVerilog w podzbiorze rozumianym przez Icarus 12, Verilatora i
    Yosysa: \code{logic}, \code{always\_comb}, \code{always\_ff},
    \code{typedef enum}, pakiety.
  \item \code{always\_comb} z przypisaniem \code{=} dla logiki kombinacyjnej,
    \code{always\_ff} z \code{<=} dla rejestrów. Bez wyjątków.
  \item Jeden zegar \code{clk}, reset \emph{synchroniczny}, aktywny w stanie
    wysokim (\code{rst}).
  \item Stałe architektury zawsze z pakietu \code{rv32i\_pkg}
    (\plik{common/rtl/rv32i\_pkg.sv}), nigdy „magiczne liczby”.
\end{itemize}
